\chapter{Literature Review and Related Work}
\section{History of Steganography}
Classical examples of steganography rely on concealing the existence of a
message inside an apparently ordinary object or communication. Anderson's
account includes messages hidden beneath writing surfaces and distinguishes
concealment from making content unreadable \cite{anderson}. The present
project uses the same broad question at a filesystem boundary: which bytes
are omitted by a routine view of a file? It does not implement the
cryptographic constructions discussed in that work.

The transition to digital media created several possible carriers, including
images, audio, documents, and system storage. Petitcolas's information-hiding
resource distinguishes applications such as watermarking, fingerprinting,
covert channels, and hidden-information detection \cite{petitcolas}.
Those applications have different objectives. A watermark may be intended
to survive ordinary processing, whereas a forensic training marker is useful
when it provides a known expected answer. Our plaintext markers were selected
for traceability, not for strong concealment against an informed examiner.

Filesystem hiding is also different from image-bit substitution. It relies
on how storage structures expose or omit byte ranges, rather than on a
perceptual model of pixels. This difference drives the evaluation: capacity
can be expressed in bytes, while correctness requires reconstructing the
appropriate stream or slack range. A perceptually unchanged document is not
by itself evidence that all associated storage bytes are unchanged.

\section{NTFS Structures Relevant to Recovery}
Microsoft describes the Master File Table as the structure holding
information needed to retrieve files. File records contain typed attributes,
and a file can require more than one record \cite{mft}. This supports a
metadata-driven approach: recovery begins by identifying an appropriate
record and attribute rather than searching blindly for a familiar content
signature. It also establishes a boundary for the project, because resolving
additional attribute records was not implemented.

The file-record header identifies the first attribute, flags, and the
relationship to a base record \cite{record}. The implementation checks those
fields before exporting content. Attribute headers distinguish resident
values from nonresident values and identify named attributes \cite{attribute}.
Our design treats this distinction as a branching point. Resident bytes
come from the validated record; nonresident bytes require mapping pairs.
Confusing the branches would produce an incorrect offset or an incomplete
recovery.

The update sequence array protects multisector structures by replacing and
later restoring sector-tail words. Microsoft's description requires checking
the sequence values before restoring saved words \cite{usa}. This is
especially important for resident content crossing a sector boundary.
Reading its apparent raw range without restoration can include the transfer
sequence word instead of the original content. The recovery implementation
therefore uses fixup-restored records, and the automated suite includes
a resident-content case crossing that boundary.

The standard-information attribute and filename attribute supply different
metadata roles \cite{standard,filename}. The project reports timestamps,
flags, and names for context, but selects actual content through
\file{$DATA}. During validation, filename metadata sizes did not always
agree with current data-stream sizes. We retained this distinction in the
reporting rules instead of treating every size field as interchangeable.

\section{Streams and Allocation Semantics}
Each NTFS stream has its own allocation, actual size, and valid-data length
\cite{streams}. Those concepts matter both for recovering ordinary content
and for avoiding disclosure of bytes the filesystem would not present as
initialized content. Microsoft documents that reads between valid-data
length and logical end return zeros \cite{fsutilfile}. Our recovery path
follows that semantic distinction: initialized bytes are read from mapped
extents, then any required logical tail is generated as zeros.

Final-cluster slack is a separate selection. Its start follows the logical
end in the last cluster that contains content. Recovering a normal file
should stop at its logical size; recovering a known slack marker should
explicitly request a bounded range outside that content. This separation
prevents a file export from silently mixing document bytes with hidden or
residual material.

\section{Related Systems and Tools}
\subsection{Sysinternals Streams}
Streams enumerates named streams associated with live files and directories;
its official documentation provides the stream-name syntax and command
options \cite{sysstreams}. It offers a straightforward way to expose ADS
on a mounted NTFS volume. The project differs in working from a raw image
and showing how the named attribute leads to the recovered bytes. Streams
was reviewed as related work, not benchmarked in this campaign.

\subsection{The Sleuth Kit}
The Sleuth Kit provides image-oriented filesystem analysis, including NTFS
attributes and alternate streams \cite{tsk}. Its \file{icat} utility can
export content by metadata address and optionally include slack
\cite{icat}. These capabilities directly overlap the broad recovery task.
Our implementation does not claim a new category of extraction capability.
Its smaller scope permits close inspection of a selected record, the
derived extents, the expected payload, and the resulting JSON report.

\subsection{Autopsy}
Autopsy provides a broader analysis environment around filesystem
capabilities, including case management and automated operations
\cite{tsk}. Such integration is valuable in a full investigation, where
organizing evidence and correlating findings matter alongside extraction.
Our command-line workflow does not replace that environment. The project
instead exposes a narrow teaching sequence that can be demonstrated without
introducing a complete case-management system.

\subsection{PhotoRec}
PhotoRec uses underlying data and recognized file formats rather than
depending entirely on filesystem metadata \cite{photorec}. This can help
when metadata is damaged, but it answers a different starting question
from selecting a named attribute in an intact record. For the deliberately
small plaintext markers in this experiment, a documented stream identity
and length provide a more direct oracle than a file-format signature.
No measured superiority over PhotoRec is asserted.

\section{Forensic Validation and Evidence Strength}
NIST's Computer Forensics Tool Testing program emphasizes specifications,
procedures, criteria, and test sets for examining tool reliability
\cite{cftt}. The project adopts that general testing discipline at a student
scale: supported outcomes are written down, failure conditions are exercised,
and test inputs are reproducible. It has not undergone CFTT certification
or an independent validation program.

SHA-256 is used as an integrity check, following the digest concept defined
by the Secure Hash Standard \cite{sha}. A recovered file's digest has its
strongest practical role when compared with a separately recorded expected
digest. Hashing an output after extraction establishes an identifier for that
output, but does not retrospectively establish its agreement with an unknown
original. For this reason, the earlier baseline-file content comparisons remained
unverified until the team confirmed the original hash matches on November 25.

\section{Gap Analysis and Project Rationale}
The reviewed tools already address ADS enumeration, raw-image examination,
or content recovery. The gap addressed here is instructional transparency
and reproducibility within the semester project. The implementation joins
raw geometry, record validation, named-stream selection, bounded slack
extraction, and explicit expected-value reporting in one inspectable path.
It avoids dependence on a mounted original disk for recovery.

That choice trades breadth for a clearer explanation of supported behavior.
It also creates identifiable weaknesses: a known-record workflow does not
find every hidden payload, and a handcrafted fixture does not reproduce
Windows allocation behavior. Those weaknesses are not shortcomings proved
against the reviewed mature tools. They are limitations of this project's
scope and evidence. The comparison in Chapter 5 therefore uses a simple,
actually executed direct-byte-read baseline, while the tools reviewed here
remain qualitative context.
